\Chapter{100}

\Verse{1}
{
    \zR{话说贾政去见节度，进去了半日，不见出来，外头议论不一。}
}
{
    \eR{[English source not present in the checked-in Bencraft/Joly files.]}
}
{
}

\Verse{2}
{
    \zR{李十儿在外也打 听不出什么事来，便想到报上的饥荒，实在也着急。}
}
{
    \eR{[English source not present in the checked-in Bencraft/Joly files.]}
}
{
}

\Verse{3}
{
    \zR{好容易听见贾政出来了，便迎 上来跟着，等不得回去，在无人处便问：“老爷进去这半天，有什么要紧的事？”}
}
{
    \eR{[English source not present in the checked-in Bencraft/Joly files.]}
}
{
}

\Verse{4}
{
    \zR{贾政笑道：“并没有事。}
}
{
    \eR{[English source not present in the checked-in Bencraft/Joly files.]}
}
{
}

\Verse{5}
{
    \zR{只为镇海总制是这位大人的亲戚，有书来嘱托照应我，所 以说了些好话。}
}
{
    \eR{[English source not present in the checked-in Bencraft/Joly files.]}
}
{
}

\Verse{6}
{
    \zR{又说：‘我们如今也是亲戚了。’”}
}
{
    \eR{[English source not present in the checked-in Bencraft/Joly files.]}
}
{
}

\Verse{7}
{
    \zR{李十儿听得，心内喜欢，不免 又壮了些胆子，便竭力怂恿贾政许这亲事。}
}
{
    \eR{[English source not present in the checked-in Bencraft/Joly files.]}
}
{
}

\Verse{8}
{
    \zR{贾政心想薛蟠的事，到底有什么挂碍，在外头信息不通，难以打点。}
}
{
    \eR{[English source not present in the checked-in Bencraft/Joly files.]}
}
{
}

\Verse{9}
{
    \zR{故回到本 任来便打发家人进京打听，顺便将总制求亲之事回明贾母，如若愿意，即将三姑娘 接到任所。}
}
{
    \eR{[English source not present in the checked-in Bencraft/Joly files.]}
}
{
}

\Verse{10}
{
    \zR{家人奉命，赶到京中回明了王夫人，便在吏部打听得贾政并无处分，惟 将署太平县的这位老爷革职。}
}
{
    \eR{[English source not present in the checked-in Bencraft/Joly files.]}
}
{
}

\Verse{11}
{
    \zR{即写了？}
}
{
    \eR{[English source not present in the checked-in Bencraft/Joly files.]}
}
{
}

\Verse{12}
{
    \zR{帖，安慰了贾政，然后住着等信。}
}
{
    \eR{[English source not present in the checked-in Bencraft/Joly files.]}
}
{
}

\Verse{13}
{
    \zR{且说薛姨妈为着薛蟠这件人命官司，各衙门内不知花了多少银钱，才定了误杀 具题。}
}
{
    \eR{[English source not present in the checked-in Bencraft/Joly files.]}
}
{
}

\Verse{14}
{
    \zR{原打量将当铺折变给人，备银赎罪，不想刑部驳审，又托人花了好些钱，总 不中用，依旧定了个死罪，监着守候秋天大审。}
}
{
    \eR{[English source not present in the checked-in Bencraft/Joly files.]}
}
{
}

\Verse{15}
{
    \zR{薛姨妈又气又疼，日夜啼哭。}
}
{
    \eR{[English source not present in the checked-in Bencraft/Joly files.]}
}
{
}

\Verse{16}
{
    \zR{宝钗 虽时常过来劝解，说是：“哥哥本来没造化。}
}
{
    \eR{[English source not present in the checked-in Bencraft/Joly files.]}
}
{
}

\Verse{17}
{
    \zR{承受了祖父这些家业，就该安安顿顿 的守着过日子。}
}
{
    \eR{[English source not present in the checked-in Bencraft/Joly files.]}
}
{
}

\Verse{18}
{
    \zR{在南边已经闹的不像样，便是香菱那件事情就了不得，因为仗着亲 戚们的势力，花了些银钱，这算白打死了一个公子。}
}
{
    \eR{[English source not present in the checked-in Bencraft/Joly files.]}
}
{
}

\Verse{19}
{
    \zR{哥哥就该改过，做起正经人来， 也该奉养母亲才是，不想进了京仍是这样。}
}
{
    \eR{[English source not present in the checked-in Bencraft/Joly files.]}
}
{
}

\Verse{20}
{
    \zR{妈妈为他不知受了多少气，哭掉了多少 眼泪。}
}
{
    \eR{[English source not present in the checked-in Bencraft/Joly files.]}
}
{
}

\Verse{21}
{
    \zR{给他娶了亲，原想大家安安逸逸的过日子，不想命该如此，偏偏娶的嫂子又 是一个不安静的，所以哥哥躲出门去。}
}
{
    \eR{[English source not present in the checked-in Bencraft/Joly files.]}
}
{
}

\Verse{22}
{
    \zR{真正俗语说的，‘冤家路儿狭’，不多几天 就闹出人命来了!妈妈和二哥哥也算不得不尽心的了：花了银钱不算，自己还求三 拜四的谋干。}
}
{
    \eR{[English source not present in the checked-in Bencraft/Joly files.]}
}
{
}

\Verse{23}
{
    \zR{无奈命里应该，也算自作自受。}
}
{
    \eR{[English source not present in the checked-in Bencraft/Joly files.]}
}
{
}

\Verse{24}
{
    \zR{大凡养儿女是为着老来有靠，便是小 户人家，还要挣一碗饭养活母亲，那里有将现成的闹光了，反害的老人家哭的死去 活来的？}
}
{
    \eR{[English source not present in the checked-in Bencraft/Joly files.]}
}
{
}

\Verse{25}
{
    \zR{不是我说，哥哥的这样行为，不是儿子，竟是个冤家对头。}
}
{
    \eR{[English source not present in the checked-in Bencraft/Joly files.]}
}
{
}

\Verse{26}
{
    \zR{妈妈再不明白， 明哭到夜，夜哭到明，又受嫂子的气。}
}
{
    \eR{[English source not present in the checked-in Bencraft/Joly files.]}
}
{
}

\Verse{27}
{
    \zR{我呢，又不能常在这里劝解。}
}
{
    \eR{[English source not present in the checked-in Bencraft/Joly files.]}
}
{
}

\Verse{28}
{
    \zR{我看见妈妈这 样，那里放得下心!他虽说是傻，也不肯叫我回去。}
}
{
    \eR{[English source not present in the checked-in Bencraft/Joly files.]}
}
{
}

\Verse{29}
{
    \zR{前儿老爷打发人回来说，看见 京报，唬的了不得，所以才叫人来打点的。}
}
{
    \eR{[English source not present in the checked-in Bencraft/Joly files.]}
}
{
}

\Verse{30}
{
    \zR{我想哥哥闹了事，担心的人也不少。}
}
{
    \eR{[English source not present in the checked-in Bencraft/Joly files.]}
}
{
}

\Verse{31}
{
    \zR{幸 亏我还是在跟前的一样，若是离乡调远，听见了这个信，只怕我想妈妈也就想杀了。}
}
{
    \eR{[English source not present in the checked-in Bencraft/Joly files.]}
}
{
}

\Verse{32}
{
    \zR{我求妈妈暂且养养神，趁哥哥的活口现在，问问各处的帐目。}
}
{
    \eR{[English source not present in the checked-in Bencraft/Joly files.]}
}
{
}

\Verse{33}
{
    \zR{人家该咱们的，咱们 该人家的，亦该请个旧伙计来算一算，看看还有几个钱没有。”}
}
{
    \eR{[English source not present in the checked-in Bencraft/Joly files.]}
}
{
}

\Verse{34}
{
    \zR{薛姨妈哭着说道： “这几天为闹你哥哥的事，你来了，不是你劝我，就是我告诉你衙门的事。}
}
{
    \eR{[English source not present in the checked-in Bencraft/Joly files.]}
}
{
}

\Verse{35}
{
    \zR{你还不 知道：京里官商的名字已经退了，两个当铺已经给了人家，银子早拿来使完了。}
}
{
    \eR{[English source not present in the checked-in Bencraft/Joly files.]}
}
{
}

\Verse{36}
{
    \zR{还 有一个当铺，管事的逃了，亏空了好几千两银子，也夹在里头打官司。}
}
{
    \eR{[English source not present in the checked-in Bencraft/Joly files.]}
}
{
}

\Verse{37}
{
    \zR{你二哥哥天 天在外头要账，料着京里的账已经去了几万银子，只好拿南边公分里银子和住房折 变才够。}
}
{
    \eR{[English source not present in the checked-in Bencraft/Joly files.]}
}
{
}

\Verse{38}
{
    \zR{前两天还听见一个荒信，说是南边的公分当铺也因为折了本儿收了。}
}
{
    \eR{[English source not present in the checked-in Bencraft/Joly files.]}
}
{
}

\Verse{39}
{
    \zR{要是 这么着，你娘的命可就活不成了！”}
}
{
    \eR{[English source not present in the checked-in Bencraft/Joly files.]}
}
{
}

\Verse{40}
{
    \zR{说着，又大哭起来。}
}
{
    \eR{[English source not present in the checked-in Bencraft/Joly files.]}
}
{
}

\Verse{41}
{
    \zR{宝钗也哭着劝道：“银钱 的事，妈妈操心也不中用，还有二哥哥给我们料理。}
}
{
    \eR{[English source not present in the checked-in Bencraft/Joly files.]}
}
{
}

\Verse{42}
{
    \zR{单可恨这些伙计们，见咱们的 势头儿败了，各自奔各自的去也罢了，我还听见说带着人家来挤我们的讹头。}
}
{
    \eR{[English source not present in the checked-in Bencraft/Joly files.]}
}
{
}

\Verse{43}
{
    \zR{可见 我哥哥活了这么大，交的人总不过是些个酒肉弟兄，急难中是一个没有的。}
}
{
    \eR{[English source not present in the checked-in Bencraft/Joly files.]}
}
{
}

\Verse{44}
{
    \zR{妈妈要 是疼我，听我的话：有年纪的人自己保重些。}
}
{
    \eR{[English source not present in the checked-in Bencraft/Joly files.]}
}
{
}

\Verse{45}
{
    \zR{妈妈这一辈子，想来还不至挨冻受饿。}
}
{
    \eR{[English source not present in the checked-in Bencraft/Joly files.]}
}
{
}

\Verse{46}
{
    \zR{家里这点子衣裳家伙，只好任凭嫂子去，那是没法儿的了。}
}
{
    \eR{[English source not present in the checked-in Bencraft/Joly files.]}
}
{
}

\Verse{47}
{
    \zR{所有的家人老婆们，瞧 他们也没心在这里了，该去的叫他们去。}
}
{
    \eR{[English source not present in the checked-in Bencraft/Joly files.]}
}
{
}

\Verse{48}
{
    \zR{只可怜香菱苦了一辈子，只好跟着妈妈。}
}
{
    \eR{[English source not present in the checked-in Bencraft/Joly files.]}
}
{
}

\Verse{49}
{
    \zR{实在短什么，我要是有的，还可以拿些个来，料我们那个也没有不依的。}
}
{
    \eR{[English source not present in the checked-in Bencraft/Joly files.]}
}
{
}

\Verse{50}
{
    \zR{就是袭姑 娘也是心术正道的，他听见咱们家的事，他倒提起妈妈来就哭。}
}
{
    \eR{[English source not present in the checked-in Bencraft/Joly files.]}
}
{
}

\Verse{51}
{
    \zR{我们那一个还打量 没事的，所以不大着急，要听见了，也是要唬个半死儿的。”}
}
{
    \eR{[English source not present in the checked-in Bencraft/Joly files.]}
}
{
}

\Verse{52}
{
    \zR{薛姨妈不等说完，便 说：“好姑娘，你可别告诉他。}
}
{
    \eR{[English source not present in the checked-in Bencraft/Joly files.]}
}
{
}

\Verse{53}
{
    \zR{他为一个林姑娘几乎没要了命，如今才好了些。}
}
{
    \eR{[English source not present in the checked-in Bencraft/Joly files.]}
}
{
}

\Verse{54}
{
    \zR{要 是他急出个原故来，不但你添一层烦恼，我越发没了依靠了。”}
}
{
    \eR{[English source not present in the checked-in Bencraft/Joly files.]}
}
{
}

\Verse{55}
{
    \zR{宝钗道：“我也是 这么想，所以总没告诉他。”}
}
{
    \eR{[English source not present in the checked-in Bencraft/Joly files.]}
}
{
}

\Verse{56}
{
    \zR{正说着，只听见金桂跑来外间屋里哭喊道：“我的命是不要的了!男人呢，已 经是没有活的分儿了。}
}
{
    \eR{[English source not present in the checked-in Bencraft/Joly files.]}
}
{
}

\Verse{57}
{
    \zR{咱们如今索性闹一闹，大伙儿到法场上去拚一拚！”}
}
{
    \eR{[English source not present in the checked-in Bencraft/Joly files.]}
}
{
}

\Verse{58}
{
    \zR{说着， 便将头往隔断板上乱撞，撞的披头散发。}
}
{
    \eR{[English source not present in the checked-in Bencraft/Joly files.]}
}
{
}

\Verse{59}
{
    \zR{气的薛姨妈白瞪着两只眼，一句话也说不 出来。}
}
{
    \eR{[English source not present in the checked-in Bencraft/Joly files.]}
}
{
}

\Verse{60}
{
    \zR{还亏了宝钗嫂子长嫂子短，好一句歹一句的劝他。}
}
{
    \eR{[English source not present in the checked-in Bencraft/Joly files.]}
}
{
}

\Verse{61}
{
    \zR{金桂道：“姑奶奶，如今 你是比不得头里的了。}
}
{
    \eR{[English source not present in the checked-in Bencraft/Joly files.]}
}
{
}

\Verse{62}
{
    \zR{你两口儿好好的过日子，我是个单身人儿，要脸做什么！”}
}
{
    \eR{[English source not present in the checked-in Bencraft/Joly files.]}
}
{
}

\Verse{63}
{
    \zR{说着，就要跑到街上回娘家去。}
}
{
    \eR{[English source not present in the checked-in Bencraft/Joly files.]}
}
{
}

\Verse{64}
{
    \zR{亏了人还多，拉住了，又劝了半天方住。}
}
{
    \eR{[English source not present in the checked-in Bencraft/Joly files.]}
}
{
}

\Verse{65}
{
    \zR{把个宝琴 唬的再不敢见他。}
}
{
    \eR{[English source not present in the checked-in Bencraft/Joly files.]}
}
{
}

\Verse{66}
{
    \zR{若是薛蝌在家，他便抹粉施脂，描眉画鬓，奇情异致的打扮收拾 起来，不时打从薛蝌住房前过，或故意咳嗽一声，明知薛蝌在屋里，特问房里是谁。}
}
{
    \eR{[English source not present in the checked-in Bencraft/Joly files.]}
}
{
}

\Verse{67}
{
    \zR{有时遇见薛蝌，他便妖妖调调、娇娇痴痴的问寒问暖，忽喜忽嗔。}
}
{
    \eR{[English source not present in the checked-in Bencraft/Joly files.]}
}
{
}

\Verse{68}
{
    \zR{丫头们看见都连 忙躲开，他自己也不觉得，只是一心一意要弄的薛蝌感情时，好行宝蟾之计。}
}
{
    \eR{[English source not present in the checked-in Bencraft/Joly files.]}
}
{
}

\Verse{69}
{
    \zR{那薛 蝌却只躲着，有时遇见也不敢不周旋他，倒是怕他撒泼放刁的意思。}
}
{
    \eR{[English source not present in the checked-in Bencraft/Joly files.]}
}
{
}

\Verse{70}
{
    \zR{更加金桂一则 为色迷心，越瞧越爱，越想越幻，那里还看的出薛蝌的真假来？}
}
{
    \eR{[English source not present in the checked-in Bencraft/Joly files.]}
}
{
}

\Verse{71}
{
    \zR{只有一宗，他见薛 蝌有什么东西都是托香菱收着，衣服缝洗也是香菱，两个人偶然说话，他来了，急 忙散开：一发动了一个“醋”字。}
}
{
    \eR{[English source not present in the checked-in Bencraft/Joly files.]}
}
{
}

\Verse{72}
{
    \zR{欲待发作薛蝌，却是舍不得，只得将一腔隐恨都 搁在香菱身上。}
}
{
    \eR{[English source not present in the checked-in Bencraft/Joly files.]}
}
{
}

\Verse{73}
{
    \zR{却又恐怕闹了香菱得罪了薛蝌，倒弄的隐忍不发。}
}
{
    \eR{[English source not present in the checked-in Bencraft/Joly files.]}
}
{
}

\Verse{74}
{
    \zR{一日，宝蟾走来，笑嘻嘻的向金桂道：“奶奶，看见了二爷没有？”}
}
{
    \eR{[English source not present in the checked-in Bencraft/Joly files.]}
}
{
}

\Verse{75}
{
    \zR{金桂道： “没有。”}
}
{
    \eR{[English source not present in the checked-in Bencraft/Joly files.]}
}
{
}

\Verse{76}
{
    \zR{宝蟾笑道：“我说二爷的那种假正经是信不得的。}
}
{
    \eR{[English source not present in the checked-in Bencraft/Joly files.]}
}
{
}

\Verse{77}
{
    \zR{咱们前儿送了酒去， 他说不会喝，刚才我见他到太太那屋里去，脸上红扑扑儿的一脸酒气。}
}
{
    \eR{[English source not present in the checked-in Bencraft/Joly files.]}
}
{
}

\Verse{78}
{
    \zR{奶奶不信， 回来只在咱们院子门口儿等他。}
}
{
    \eR{[English source not present in the checked-in Bencraft/Joly files.]}
}
{
}

\Verse{79}
{
    \zR{他打那边过来，奶奶叫住他问问，看他说什么。”}
}
{
    \eR{[English source not present in the checked-in Bencraft/Joly files.]}
}
{
}

\Verse{80}
{
    \zR{金桂听了，一心的恼意，便道：“他那里就出来了呢。}
}
{
    \eR{[English source not present in the checked-in Bencraft/Joly files.]}
}
{
}

\Verse{81}
{
    \zR{他既无情义，问他作什么？”}
}
{
    \eR{[English source not present in the checked-in Bencraft/Joly files.]}
}
{
}

\Verse{82}
{
    \zR{宝蟾道：“奶奶又迂了。}
}
{
    \eR{[English source not present in the checked-in Bencraft/Joly files.]}
}
{
}

\Verse{83}
{
    \zR{他好说，咱们也好说；他不好说，咱们再另打主意。”}
}
{
    \eR{[English source not present in the checked-in Bencraft/Joly files.]}
}
{
}

\Verse{84}
{
    \zR{金 桂听着有理，因叫宝蟾：“瞧着他，看他出去了。”}
}
{
    \eR{[English source not present in the checked-in Bencraft/Joly files.]}
}
{
}

\Verse{85}
{
    \zR{宝蟾答应着出来，金桂却去打 开镜奁，又照了一照，把嘴唇儿又抹了一抹。}
}
{
    \eR{[English source not present in the checked-in Bencraft/Joly files.]}
}
{
}

\Verse{86}
{
    \zR{然后拿一条洒花绢子，才要出来，又 像忘了什么的，心里倒不知怎么是好了。}
}
{
    \eR{[English source not present in the checked-in Bencraft/Joly files.]}
}
{
}

\Verse{87}
{
    \zR{只听宝蟾外面说道：“二爷今日高兴啊。}
}
{
    \eR{[English source not present in the checked-in Bencraft/Joly files.]}
}
{
}

\Verse{88}
{
    \zR{那里喝了酒来了？”}
}
{
    \eR{[English source not present in the checked-in Bencraft/Joly files.]}
}
{
}

\Verse{89}
{
    \zR{金桂听了，明知是叫他出来的意思，连忙掀起帘子出来。}
}
{
    \eR{[English source not present in the checked-in Bencraft/Joly files.]}
}
{
}

\Verse{90}
{
    \zR{只见 薛蝌和宝蟾说道：“今日是张大爷的好日子，所以被他们强不过，吃了半钟。}
}
{
    \eR{[English source not present in the checked-in Bencraft/Joly files.]}
}
{
}

\Verse{91}
{
    \zR{到这 时候脸还发烧呢。”}
}
{
    \eR{[English source not present in the checked-in Bencraft/Joly files.]}
}
{
}

\Verse{92}
{
    \zR{一句话没说完，金桂早接口道：“自然人家外人的酒，比咱们 自己家里的酒是有趣儿的。”}
}
{
    \eR{[English source not present in the checked-in Bencraft/Joly files.]}
}
{
}

\Verse{93}
{
    \zR{薛蝌被他拿话一激，脸越红了，连忙走过来陪笑道： “嫂子说那里的话？”}
}
{
    \eR{[English source not present in the checked-in Bencraft/Joly files.]}
}
{
}

\Verse{94}
{
    \zR{宝蟾见他二人交谈，便躲到屋里去了。}
}
{
    \eR{[English source not present in the checked-in Bencraft/Joly files.]}
}
{
}

\Verse{95}
{
    \zR{这金桂初时原要假意 发作薛蝌两句，无奈一见他两颊微红，双眸带涩，别有一种谨愿可怜之意，早把自
        己那骄悍之气，感化到爪洼国去了，因笑说道：“这么说，你的酒是硬强着才肯喝 的呢。”}
}
{
    \eR{[English source not present in the checked-in Bencraft/Joly files.]}
}
{
}

\Verse{96}
{
    \zR{薛蝌道：“我那里喝得来？”}
}
{
    \eR{[English source not present in the checked-in Bencraft/Joly files.]}
}
{
}

\Verse{97}
{
    \zR{金桂道：“不喝也好，强如像你哥哥喝出乱 子来，明儿娶了你们奶奶儿，像我这样守活寡受孤单呢！”}
}
{
    \eR{[English source not present in the checked-in Bencraft/Joly files.]}
}
{
}

\Verse{98}
{
    \zR{说到这里，两个眼已经 乜斜了，两腮上也觉红晕了。}
}
{
    \eR{[English source not present in the checked-in Bencraft/Joly files.]}
}
{
}

\Verse{99}
{
    \zR{薛蝌见这话越发邪僻了，打算着要走。}
}
{
    \eR{[English source not present in the checked-in Bencraft/Joly files.]}
}
{
}

\Verse{100}
{
    \zR{金桂也看出来 了，那里容得，早已走过来一把拉住。}
}
{
    \eR{[English source not present in the checked-in Bencraft/Joly files.]}
}
{
}

\Verse{101}
{
    \zR{薛蝌急了道：“嫂子放尊重些。”}
}
{
    \eR{[English source not present in the checked-in Bencraft/Joly files.]}
}
{
}

\Verse{102}
{
    \zR{说着浑身 乱颤。}
}
{
    \eR{[English source not present in the checked-in Bencraft/Joly files.]}
}
{
}

\Verse{103}
{
    \zR{金桂索性老着脸道：“你只管进来，我和你说一句要紧的话。”}
}
{
    \eR{[English source not present in the checked-in Bencraft/Joly files.]}
}
{
}

\Verse{104}
{
    \zR{正闹着，忽听背后一个人叫道：“奶奶!香菱来了。”}
}
{
    \eR{[English source not present in the checked-in Bencraft/Joly files.]}
}
{
}

\Verse{105}
{
    \zR{把金桂唬了一跳。}
}
{
    \eR{[English source not present in the checked-in Bencraft/Joly files.]}
}
{
}

\Verse{106}
{
    \zR{回头 瞧时，却是宝蟾掀着帘子看他二人的光景，一抬头见香菱从那边来了，赶忙知会金 桂。}
}
{
    \eR{[English source not present in the checked-in Bencraft/Joly files.]}
}
{
}

\Verse{107}
{
    \zR{金桂这一惊不小，手已松了。}
}
{
    \eR{[English source not present in the checked-in Bencraft/Joly files.]}
}
{
}

\Verse{108}
{
    \zR{薛蝌得便脱身跑了。}
}
{
    \eR{[English source not present in the checked-in Bencraft/Joly files.]}
}
{
}

\Verse{109}
{
    \zR{那香菱正走着，原不理会， 忽听宝蟾一嚷，才瞧见金桂在那里拉住薛蝌，往里死拽。}
}
{
    \eR{[English source not present in the checked-in Bencraft/Joly files.]}
}
{
}

\Verse{110}
{
    \zR{香菱却唬的心头乱跳，自 己连忙转身回去。}
}
{
    \eR{[English source not present in the checked-in Bencraft/Joly files.]}
}
{
}

\Verse{111}
{
    \zR{这里金桂早已连吓带气，呆呆的瞅着薛蝌去了，怔了半天，恨了 一声，自己扫兴归房。}
}
{
    \eR{[English source not present in the checked-in Bencraft/Joly files.]}
}
{
}

\Verse{112}
{
    \zR{从此把香菱恨入骨髓。}
}
{
    \eR{[English source not present in the checked-in Bencraft/Joly files.]}
}
{
}

\Verse{113}
{
    \zR{那香菱本是要到宝琴那里，刚走出腰 门，看见这般，吓回去了。}
}
{
    \eR{[English source not present in the checked-in Bencraft/Joly files.]}
}
{
}

\Verse{114}
{
    \zR{是日，宝钗在贾母屋里，听得王夫人告诉老太太要聘探春一事。}
}
{
    \eR{[English source not present in the checked-in Bencraft/Joly files.]}
}
{
}

\Verse{115}
{
    \zR{贾母说道：“既 是同乡的人，很好。}
}
{
    \eR{[English source not present in the checked-in Bencraft/Joly files.]}
}
{
}

\Verse{116}
{
    \zR{只是听见说那孩子到过我们家里，怎么你老爷没有提起？”}
}
{
    \eR{[English source not present in the checked-in Bencraft/Joly files.]}
}
{
}

\Verse{117}
{
    \zR{王 夫人道：“连我们也不知道。”}
}
{
    \eR{[English source not present in the checked-in Bencraft/Joly files.]}
}
{
}

\Verse{118}
{
    \zR{贾母道：“好是好，但只道儿太远。}
}
{
    \eR{[English source not present in the checked-in Bencraft/Joly files.]}
}
{
}

\Verse{119}
{
    \zR{虽然老爷在那 里，倘或将来老爷调任，可不是我们孩子太单了吗？”}
}
{
    \eR{[English source not present in the checked-in Bencraft/Joly files.]}
}
{
}

\Verse{120}
{
    \zR{王夫人道：“两家都是做官 的，也是拿不定。}
}
{
    \eR{[English source not present in the checked-in Bencraft/Joly files.]}
}
{
}

\Verse{121}
{
    \zR{或者那边还调进来，即不然，终有个叶落归根。}
}
{
    \eR{[English source not present in the checked-in Bencraft/Joly files.]}
}
{
}

\Verse{122}
{
    \zR{况且老爷既在那 里做官，上司已经说了，好意思不给么？}
}
{
    \eR{[English source not present in the checked-in Bencraft/Joly files.]}
}
{
}

\Verse{123}
{
    \zR{想来老爷的主意定了，只是不敢做主，故 遣人来回老太太的。”}
}
{
    \eR{[English source not present in the checked-in Bencraft/Joly files.]}
}
{
}

\Verse{124}
{
    \zR{贾母道：“你们愿意更好，但是三丫头这一去了，不知三年 两年那边可能回家？}
}
{
    \eR{[English source not present in the checked-in Bencraft/Joly files.]}
}
{
}

\Verse{125}
{
    \zR{若再迟了，恐怕我赶不上再见他一面了。”}
}
{
    \eR{[English source not present in the checked-in Bencraft/Joly files.]}
}
{
}

\Verse{126}
{
    \zR{说着掉下泪来。}
}
{
    \eR{[English source not present in the checked-in Bencraft/Joly files.]}
}
{
}

\Verse{127}
{
    \zR{王 夫人道：“孩子们大了，少不得总要给人家的。}
}
{
    \eR{[English source not present in the checked-in Bencraft/Joly files.]}
}
{
}

\Verse{128}
{
    \zR{就是本乡本土的人，除非不做官还 使得，要是做官的，谁保的住总在一处？}
}
{
    \eR{[English source not present in the checked-in Bencraft/Joly files.]}
}
{
}

\Verse{129}
{
    \zR{只要孩子们有造化就好。}
}
{
    \eR{[English source not present in the checked-in Bencraft/Joly files.]}
}
{
}

\Verse{130}
{
    \zR{譬如迎姑娘倒配 的近呢，偏时常听见他和女婿打闹，甚至于不给饭吃。}
}
{
    \eR{[English source not present in the checked-in Bencraft/Joly files.]}
}
{
}

\Verse{131}
{
    \zR{就是我们送了东西去，他也 摸不着。}
}
{
    \eR{[English source not present in the checked-in Bencraft/Joly files.]}
}
{
}

\Verse{132}
{
    \zR{近来听见益发不好了，也不放他回来。}
}
{
    \eR{[English source not present in the checked-in Bencraft/Joly files.]}
}
{
}

\Verse{133}
{
    \zR{两口子拌起来，就说咱们使了他家 的银钱，可怜这孩子总不得个出头的日子。}
}
{
    \eR{[English source not present in the checked-in Bencraft/Joly files.]}
}
{
}

\Verse{134}
{
    \zR{前儿我掂记他，打发人去瞧他，迎丫头 藏在耳房里，不肯出来。}
}
{
    \eR{[English source not present in the checked-in Bencraft/Joly files.]}
}
{
}

\Verse{135}
{
    \zR{老婆们必要进去，看见我们姑娘这样冷天还穿着几件旧衣 裳。}
}
{
    \eR{[English source not present in the checked-in Bencraft/Joly files.]}
}
{
}

\Verse{136}
{
    \zR{他一包眼泪的告诉老婆们说：‘回去别说我这么苦，这也是我命里所招!也不 用送什么衣裳东西来，不但摸不着，反要添一顿打，说是我告诉的。’}
}
{
    \eR{[English source not present in the checked-in Bencraft/Joly files.]}
}
{
}

\Verse{137}
{
    \zR{老太太想想， 这倒是近处眼见的，若不好，更难受。}
}
{
    \eR{[English source not present in the checked-in Bencraft/Joly files.]}
}
{
}

\Verse{138}
{
    \zR{倒亏了大太太也不理会他，大老爷也不出个 头。}
}
{
    \eR{[English source not present in the checked-in Bencraft/Joly files.]}
}
{
}

\Verse{139}
{
    \zR{如今迎姑娘实在比我们三等使唤的丫头还不及。}
}
{
    \eR{[English source not present in the checked-in Bencraft/Joly files.]}
}
{
}

\Verse{140}
{
    \zR{我想探丫头虽不是我养的，老 爷既看见过女婿，定然是好才许的。}
}
{
    \eR{[English source not present in the checked-in Bencraft/Joly files.]}
}
{
}

\Verse{141}
{
    \zR{只请老太太示下，择个好日子，多派几个人送 到他老爷任上，该怎么着，老爷也不肯将就。”}
}
{
    \eR{[English source not present in the checked-in Bencraft/Joly files.]}
}
{
}

\Verse{142}
{
    \zR{贾母道：“有他老子作主，你就料 理妥当，拣个长行的日子送去，也就定了一件事。”}
}
{
    \eR{[English source not present in the checked-in Bencraft/Joly files.]}
}
{
}

\Verse{143}
{
    \zR{王夫人答应着“是”。}
}
{
    \eR{[English source not present in the checked-in Bencraft/Joly files.]}
}
{
}

\Verse{144}
{
    \zR{宝钗听 的明白，也不敢则声，只是心里叫苦：“我们家的姑娘们就算他是个尖儿。}
}
{
    \eR{[English source not present in the checked-in Bencraft/Joly files.]}
}
{
}

\Verse{145}
{
    \zR{如今又 要远嫁，眼看着这里的人一天少似一天了。”}
}
{
    \eR{[English source not present in the checked-in Bencraft/Joly files.]}
}
{
}

\Verse{146}
{
    \zR{见王夫人起身告辞出去，他也送出来 了。}
}
{
    \eR{[English source not present in the checked-in Bencraft/Joly files.]}
}
{
}

\Verse{147}
{
    \zR{一径回到自己房中，并不与宝玉说知，见袭人独自一个做活，便将听见的话说 了。}
}
{
    \eR{[English source not present in the checked-in Bencraft/Joly files.]}
}
{
}

\Verse{148}
{
    \zR{袭人也很不受用。}
}
{
    \eR{[English source not present in the checked-in Bencraft/Joly files.]}
}
{
}

\Verse{149}
{
    \zR{却说赵姨娘听见探春这事，反喜欢起来，心里说道：“我这个丫头在家忒瞧不 起我，我何从还是个娘？}
}
{
    \eR{[English source not present in the checked-in Bencraft/Joly files.]}
}
{
}

\Verse{150}
{
    \zR{比他的丫头还不济。}
}
{
    \eR{[English source not present in the checked-in Bencraft/Joly files.]}
}
{
}

\Verse{151}
{
    \zR{况且？}
}
{
    \eR{[English source not present in the checked-in Bencraft/Joly files.]}
}
{
}

\Verse{152}
{
    \zR{上水，护着别人。}
}
{
    \eR{[English source not present in the checked-in Bencraft/Joly files.]}
}
{
}

\Verse{153}
{
    \zR{他挡在头里， 连环儿也不得出头。}
}
{
    \eR{[English source not present in the checked-in Bencraft/Joly files.]}
}
{
}

\Verse{154}
{
    \zR{如今老爷接了去，我倒干净。}
}
{
    \eR{[English source not present in the checked-in Bencraft/Joly files.]}
}
{
}

\Verse{155}
{
    \zR{想要他孝敬我不能够了，只愿意 他像迎丫头似的，我也称称愿。”}
}
{
    \eR{[English source not present in the checked-in Bencraft/Joly files.]}
}
{
}

\Verse{156}
{
    \zR{一面想着，一面跑到探春那边与他道喜，说：“姑 娘，你是要高飞的人了。}
}
{
    \eR{[English source not present in the checked-in Bencraft/Joly files.]}
}
{
}

\Verse{157}
{
    \zR{到了姑爷那边自然比家里还好，想来你也是愿意的。}
}
{
    \eR{[English source not present in the checked-in Bencraft/Joly files.]}
}
{
}

\Verse{158}
{
    \zR{就是 养了你一场，并没有借你的光儿。}
}
{
    \eR{[English source not present in the checked-in Bencraft/Joly files.]}
}
{
}

\Verse{159}
{
    \zR{就是我有七分不好，也有三分的好，也别说一去 了把我搁在脑杓子后头。”}
}
{
    \eR{[English source not present in the checked-in Bencraft/Joly files.]}
}
{
}

\Verse{160}
{
    \zR{探春听着毫无道理，只低头作活，一句也不言语。}
}
{
    \eR{[English source not present in the checked-in Bencraft/Joly files.]}
}
{
}

\Verse{161}
{
    \zR{赵姨 娘见他不理，气忿忿的自己去了。}
}
{
    \eR{[English source not present in the checked-in Bencraft/Joly files.]}
}
{
}

\Verse{162}
{
    \zR{这里探春又气又笑又伤心，也不过自己掉泪而已。}
}
{
    \eR{[English source not present in the checked-in Bencraft/Joly files.]}
}
{
}

\Verse{163}
{
    \zR{坐了一回，闷闷的走到宝玉 这边来。}
}
{
    \eR{[English source not present in the checked-in Bencraft/Joly files.]}
}
{
}

\Verse{164}
{
    \zR{宝玉因问道：“三妹妹，我听见林妹妹死的时候，你在那里来着。}
}
{
    \eR{[English source not present in the checked-in Bencraft/Joly files.]}
}
{
}

\Verse{165}
{
    \zR{我还听 见说：林妹妹死的时候，远远的有音乐之声。}
}
{
    \eR{[English source not present in the checked-in Bencraft/Joly files.]}
}
{
}

\Verse{166}
{
    \zR{或者他是有来历的，也未可知。”}
}
{
    \eR{[English source not present in the checked-in Bencraft/Joly files.]}
}
{
}

\Verse{167}
{
    \zR{探 春笑道：“那是你心里想着罢了。}
}
{
    \eR{[English source not present in the checked-in Bencraft/Joly files.]}
}
{
}

\Verse{168}
{
    \zR{但只那夜却怪，不像人家鼓乐的声儿，你的话或 者也是。”}
}
{
    \eR{[English source not present in the checked-in Bencraft/Joly files.]}
}
{
}

\Verse{169}
{
    \zR{宝玉听了，更以为实。}
}
{
    \eR{[English source not present in the checked-in Bencraft/Joly files.]}
}
{
}

\Verse{170}
{
    \zR{又想前日自己神魂飘荡之时，曾见一人，说是黛 玉生不同人，死不同鬼，必是那里的仙子临凡。}
}
{
    \eR{[English source not present in the checked-in Bencraft/Joly files.]}
}
{
}

\Verse{171}
{
    \zR{又想起那年唱戏做的嫦娥，飘飘艳 艳，何等风致。}
}
{
    \eR{[English source not present in the checked-in Bencraft/Joly files.]}
}
{
}

\Verse{172}
{
    \zR{过了一回探春去了，因必要紫鹃过来，立刻回了贾母去叫他。}
}
{
    \eR{[English source not present in the checked-in Bencraft/Joly files.]}
}
{
}

\Verse{173}
{
    \zR{无奈 紫鹃心里不愿意，虽经贾母王夫人派了过来，自己没法，却是在宝玉跟前，不是嗳 声就是叹气的。}
}
{
    \eR{[English source not present in the checked-in Bencraft/Joly files.]}
}
{
}

\Verse{174}
{
    \zR{宝玉背地里拉着他，低声下气要问黛玉的话，紫鹃从没好话回答。}
}
{
    \eR{[English source not present in the checked-in Bencraft/Joly files.]}
}
{
}

\Verse{175}
{
    \zR{宝钗倒背地里夸他有忠心，并不嗔怪他。}
}
{
    \eR{[English source not present in the checked-in Bencraft/Joly files.]}
}
{
}

\Verse{176}
{
    \zR{那雪雁虽是宝玉娶亲这夜出过力的，宝玉 见他心地不甚明白，便回了贾母王夫人，将他配了一个小厮，各自过活去了。}
}
{
    \eR{[English source not present in the checked-in Bencraft/Joly files.]}
}
{
}

\Verse{177}
{
    \zR{王奶 妈，养着他将来好送黛玉的灵柩回南。}
}
{
    \eR{[English source not present in the checked-in Bencraft/Joly files.]}
}
{
}

\Verse{178}
{
    \zR{鹦哥等小丫头，仍旧伏侍老太太。}
}
{
    \eR{[English source not present in the checked-in Bencraft/Joly files.]}
}
{
}

\Verse{179}
{
    \zR{宝玉本想念黛玉，因此及彼，又想跟黛玉的人已经云散，更加纳闷。}
}
{
    \eR{[English source not present in the checked-in Bencraft/Joly files.]}
}
{
}

\Verse{180}
{
    \zR{闷到无可 如何，忽又想黛玉死的这样清楚，必是离凡返仙去了，反又欢喜。}
}
{
    \eR{[English source not present in the checked-in Bencraft/Joly files.]}
}
{
}

\Verse{181}
{
    \zR{忽然听见袭人和 宝钗那里讲究探春出嫁之事，宝玉听了，“啊呀”的一声，哭倒在炕上。}
}
{
    \eR{[English source not present in the checked-in Bencraft/Joly files.]}
}
{
}

\Verse{182}
{
    \zR{唬得宝钗 袭人都来扶起，说：“怎么了？”}
}
{
    \eR{[English source not present in the checked-in Bencraft/Joly files.]}
}
{
}

\Verse{183}
{
    \zR{宝玉早哭的说不出来。}
}
{
    \eR{[English source not present in the checked-in Bencraft/Joly files.]}
}
{
}

\Verse{184}
{
    \zR{定了一回子神，说道：“这 日子过不得了，我姊妹们都一个一个的散了!林妹妹是成了仙去了。}
}
{
    \eR{[English source not present in the checked-in Bencraft/Joly files.]}
}
{
}

\Verse{185}
{
    \zR{大姐姐呢，已 经死了，这也罢了，没天天在一块儿。}
}
{
    \eR{[English source not present in the checked-in Bencraft/Joly files.]}
}
{
}

\Verse{186}
{
    \zR{二姐姐碰着了一个混账不堪的东西。}
}
{
    \eR{[English source not present in the checked-in Bencraft/Joly files.]}
}
{
}

\Verse{187}
{
    \zR{三妹妹 又要远嫁，总不得见的了。}
}
{
    \eR{[English source not present in the checked-in Bencraft/Joly files.]}
}
{
}

\Verse{188}
{
    \zR{史妹妹又不知要到那里去。}
}
{
    \eR{[English source not present in the checked-in Bencraft/Joly files.]}
}
{
}

\Verse{189}
{
    \zR{薛妹妹是有了人家儿的。}
}
{
    \eR{[English source not present in the checked-in Bencraft/Joly files.]}
}
{
}

\Verse{190}
{
    \zR{这 些姐姐妹妹，难道一个都不留在家里，单留我做什么？”}
}
{
    \eR{[English source not present in the checked-in Bencraft/Joly files.]}
}
{
}

\Verse{191}
{
    \zR{袭人忙又拿话解劝。}
}
{
    \eR{[English source not present in the checked-in Bencraft/Joly files.]}
}
{
}

\Verse{192}
{
    \zR{宝钗 摆着手说：“你不用劝他，等我问他。”}
}
{
    \eR{[English source not present in the checked-in Bencraft/Joly files.]}
}
{
}

\Verse{193}
{
    \zR{因问着宝玉道：“据你的心里，要这些姐 妹都在家里陪到你老了，都不为终身的事吗？}
}
{
    \eR{[English source not present in the checked-in Bencraft/Joly files.]}
}
{
}

\Verse{194}
{
    \zR{要说别人，或者还有别的想头。}
}
{
    \eR{[English source not present in the checked-in Bencraft/Joly files.]}
}
{
}

\Verse{195}
{
    \zR{你自 己的姐姐妹妹，不用说没有远嫁的；就是有，老爷作主，你有什么法儿？}
}
{
    \eR{[English source not present in the checked-in Bencraft/Joly files.]}
}
{
}

\Verse{196}
{
    \zR{打量天下 就是你一个人爱姐姐妹妹呢？}
}
{
    \eR{[English source not present in the checked-in Bencraft/Joly files.]}
}
{
}

\Verse{197}
{
    \zR{要是都像你，就连我也不能陪着你了。}
}
{
    \eR{[English source not present in the checked-in Bencraft/Joly files.]}
}
{
}

\Verse{198}
{
    \zR{大凡人念书原 为的是明理，怎么你越念越糊涂了呢。}
}
{
    \eR{[English source not present in the checked-in Bencraft/Joly files.]}
}
{
}

\Verse{199}
{
    \zR{这么说起来，我和袭姑娘各自一边儿去，让 你把姐姐妹妹们都邀了来守着你。”}
}
{
    \eR{[English source not present in the checked-in Bencraft/Joly files.]}
}
{
}

\Verse{200}
{
    \zR{宝玉听了，两只手拉住宝钗袭人道：“我也知 道。}
}
{
    \eR{[English source not present in the checked-in Bencraft/Joly files.]}
}
{
}

\Verse{201}
{
    \zR{为什么散的这么早呢？}
}
{
    \eR{[English source not present in the checked-in Bencraft/Joly files.]}
}
{
}

\Verse{202}
{
    \zR{等我化了灰的时候再散也不迟。”}
}
{
    \eR{[English source not present in the checked-in Bencraft/Joly files.]}
}
{
}

\Verse{203}
{
    \zR{袭人掩着他的嘴道： “又胡说了。}
}
{
    \eR{[English source not present in the checked-in Bencraft/Joly files.]}
}
{
}

\Verse{204}
{
    \zR{才这两天身上好些，二奶奶才吃些饭。}
}
{
    \eR{[English source not present in the checked-in Bencraft/Joly files.]}
}
{
}

\Verse{205}
{
    \zR{你要是又闹翻了，我也不管了。”}
}
{
    \eR{[English source not present in the checked-in Bencraft/Joly files.]}
}
{
}

\Verse{206}
{
    \zR{宝玉听他两个人说话都有道理，只是心上不知道怎么着才好，只得说道：“我却明 白，但只是心里闹得慌。”}
}
{
    \eR{[English source not present in the checked-in Bencraft/Joly files.]}
}
{
}

\Verse{207}
{
    \zR{宝钗也不理他，暗叫袭人快把定心丸给他吃了，慢慢的 开导他。}
}
{
    \eR{[English source not present in the checked-in Bencraft/Joly files.]}
}
{
}

\Verse{208}
{
    \zR{袭人便欲告诉探春，说临行不必来辞。}
}
{
    \eR{[English source not present in the checked-in Bencraft/Joly files.]}
}
{
}

\Verse{209}
{
    \zR{宝钗道：“这怕什么？}
}
{
    \eR{[English source not present in the checked-in Bencraft/Joly files.]}
}
{
}

\Verse{210}
{
    \zR{等消停几日， 他心里明白了，还要叫他们多说句话儿呢。}
}
{
    \eR{[English source not present in the checked-in Bencraft/Joly files.]}
}
{
}

\Verse{211}
{
    \zR{况且三姑娘是极明白的人，不像那些假 惺惺的人，少不得有一番箴谏，他以后就不是这样了。”}
}
{
    \eR{[English source not present in the checked-in Bencraft/Joly files.]}
}
{
}

\Verse{212}
{
    \zR{正说着，贾母那边打发过 鸳鸯来说：“知道宝玉旧病又发，叫袭人劝说安慰，叫他不用胡思乱想。”}
}
{
    \eR{[English source not present in the checked-in Bencraft/Joly files.]}
}
{
}

\Verse{213}
{
    \zR{袭人等 应了。}
}
{
    \eR{[English source not present in the checked-in Bencraft/Joly files.]}
}
{
}

\Verse{214}
{
    \zR{鸳鸯坐了一会子去了。}
}
{
    \eR{[English source not present in the checked-in Bencraft/Joly files.]}
}
{
}

\Verse{215}
{
    \zR{那贾母又想起探春远行，虽不全备妆奁，其一应动用之物俱该预备，便把凤姐 叫来，将老爷的主意告诉了一遍，叫他料理去。}
}
{
    \eR{[English source not present in the checked-in Bencraft/Joly files.]}
}
{
}

\Verse{216}
{
    \zR{凤姐答应。}
}
{
    \eR{[English source not present in the checked-in Bencraft/Joly files.]}
}
{
}

\Verse{217}
{
    \zR{不知怎么办理，下回分解。}
}
{
    \eR{[English source not present in the checked-in Bencraft/Joly files.]}
}
{
}

\Verse{218}
{
    \zR{(本章完)}
}
{
    \eR{[English source not present in the checked-in Bencraft/Joly files.]}
}
{
}
